\section{问题三：风格条件生成与跨风格迁移}

\noindent\textbf{待 M6 补。}

本章只登记题面要求与拟采用路线，不报告任何生成结果。问题三的实验、图表与数值留待后续里程碑完成。

\subsection{题面要求}

题面要求以啤酒风格为条件，学习一个能够生成化合物谱与对应感官画像的模型，并考察生成结果在不同风格之间的迁移能力。生成结果需要满足三类约束：其一，化合物谱的维度、量纲与训练数据一致，不能生成数据中不存在的化合物；其二，感官画像要与给定风格的目标一致，不能只生成“平均啤酒”；其三，跨风格迁移要能保留源风格的主要特征，同时把目标风格的关键差异迁移过去。

与问题一、问题二相比，这一问的难点从“判别”转为“生成”。判别模型只需要给出一个条件均值或一个排序，生成模型却要刻画整个条件分布；如果只用均方误差训练，生成结果会趋向条件均值，失去样本多样性。另一方面，风格标签的样本量不均衡，少数风格可能只有很少的样本，生成器容易记住多数风格而忽略少数风格。题面还要求化学可解释，这意味着生成器不能只追求与感官标签的拟合，还要满足非负、量纲一致、关键化合物可追溯等约束。

\subsection{拟采用路线}

本文拟采用“风格条件编码—化合物谱生成—感官回评—化学约束”的四段式路线。

第一段是风格条件编码。把风格标签映射为条件向量，条件向量既可以直接用 one-hot 编码，也可以用风格共现结构做嵌入。对样本量较少的风格，嵌入可以共享相似风格的统计强度，避免每个风格单独一套参数。

第二段是化合物谱生成。优先尝试条件变分自编码器：编码器把真实化合物谱压到隐空间，解码器在风格条件与隐变量的共同作用下重建化合物谱；训练目标由重构误差与隐空间正则项组成。若条件变分自编码器的生成样本在化学有效性上不足，再比较条件扩散模型；后者的生成质量通常更好，但采样成本更高，且对超参数更敏感。无论选哪一种，都把问题二得到的关键化合物作为条件中的显式输入，而不是让生成器完全自己发现。

第三段是感官回评。把生成化合物谱送入问题一的随机森林预测器，得到感官画像；再把画像与目标风格的代表性画像比较。回评器在这里只作为“代理评委”，不能替代真实品评；它的作用是快速筛掉明显偏离目标风格的生成样本，并为生成器提供可微或可采样的反馈信号。

第四段是化学约束。对生成化合物谱施加非负约束、量纲约束与聚合量一致性约束；对题面点名的缺陷化合物，按已登记的阈值表检查是否超阈值；对无法从数据确认的约束，明确写成“资料未覆盖”，不用猜测补位。

\subsection{跨风格迁移与验证设计}

跨风格迁移拟采用“源风格编码 + 目标风格条件”的方式：固定源风格的隐变量或内容变量，替换目标风格的条件向量，观察生成化合物谱与感官画像的变化。为了避免把“换了一个随机种子”误读为迁移，验证时固定随机源，并对同一源样本做多次条件替换，比较迁移前后关键化合物的变化方向。

验证设计分四层。第一层是化学有效性：生成谱是否非负、量纲是否与训练数据一致、聚合量是否与分项大致协调。第二层是感官一致性：回评画像与目标风格的代表性画像在主要感官维度上是否同向。第三层是多样性：同一风格条件下生成的不同样本不能完全相同，否则说明模型只记住了条件均值。第四层是迁移性：源风格的主要特征在迁移后是否保留，目标风格的关键差异是否出现。四层都只写验证方案，具体指标与结果留待后续里程碑。

\subsection{风险与边界}

主要风险有三类。第一，风格样本不均衡导致少数风格生成质量差；应对方式是把风格分层纳入采样与评估，而不是只看总体平均。第二，生成器只拟合感官标签、忽略化学约束，产生“感官像、化学不像”的样本；应对方式是把化学约束放在生成之后、回评之前，先过滤再评价。第三，回评器本身有误差，生成器可能利用回评器的弱维“钻空子”；应对方式是对回评器的不确定维度降权，并用独立证据做交叉核对。

本章不给出任何生成结果、图表或数值。问题三的模型选择、参数、生成样本与跨风格迁移实验，均在第 M6 阶段完成并并入论文。
